% ---------------------------------------------------------------------
% Planificación y presupuesto
% ---------------------------------------------------------------------

\chapter{Planificación y presupuesto}
\label{cap:planificacion-presupuesto}

\begin{Resumen}
Este capítulo describe la planificación temporal del proyecto por fases, los principales riesgos técnicos identificados y las medidas adoptadas para mitigarlos, y el presupuesto de los componentes empleados en la plataforma.
\end{Resumen}

\section{Planificación del proyecto}
\label{sec:planificacion-proyecto}

El proyecto puede dividirse en las siguientes fases de trabajo, distribuidas en el tiempo como recoge la \autoref{fig:gantt-proyecto}:

\begin{enumerate}
\item \textbf{Análisis inicial}: definición del alcance, selección de componentes y estudio de alternativas de navegación visual.
\item \textbf{Montaje hardware}: ensamblaje del chasis, motores, controlador, alimentación, GPS, cámara y Raspberry Pi.
\item \textbf{Configuración de vuelo}: calibración del controlador, pruebas de motores, parámetros de seguridad y validación en tierra.
\item \textbf{Integración software}: configuración de MAVProxy, Tailscale, \texttt{tmux}, captura de imágenes y registro de telemetría.
\item \textbf{Captura experimental}: vuelos de prueba, generación de datasets y revisión de la calidad de los datos.
\item \textbf{Evaluación CVGL}: creación de mapa base, galería satelital, recuperación global, matching local y análisis de error.
\item \textbf{Redacción}: documentación de diseño, resultados, conclusiones y líneas futuras.
\end{enumerate}

\begin{figure}[htbp]
\centering
\includegraphics[width=\textwidth]{gantt_proyecto.png}
\caption{Planificación temporal del proyecto, desde la compra de componentes hasta la entrega de la memoria.}
\label{fig:gantt-proyecto}
\end{figure}

\section{Riesgos y medidas de mitigación}
\label{sec:riesgos-mitigacion}

La \autoref{tab:riesgos} resume los riesgos técnicos principales y las medidas previstas para reducir su impacto.

\begin{table}[htbp]
\centering
\caption{Riesgos técnicos del proyecto.}
\label{tab:riesgos}
\begin{tabularx}{\textwidth}{p{0.26\textwidth}X X}
\toprule
\textbf{Riesgo} & \textbf{Impacto} & \textbf{Mitigación} \\
\midrule
Vibraciones en cámara & Imágenes borrosas y peor matching visual. & Montaje antivibración, revisión de hélices y vuelos de prueba. \\
Pérdida de conexión remota & Interrupción de monitorización. & Uso de \texttt{tmux} y registro local en la Raspberry Pi. \\
Sobrecarga computacional & Imposibilidad de CVGL en tiempo real a bordo. & Evaluación fuera de bordo y propuesta de optimización futura. \\
Fallo de telemetría & Datos incompletos para evaluación. & Registro local, pruebas de MAVLink en tierra y comprobación previa al vuelo. \\
Autonomía insuficiente & Menor tiempo de captura. & Vuelos cortos e incrementales, control de carga y revisión de peso. \\
\bottomrule
\end{tabularx}
\end{table}

\section{Presupuesto}
\label{sec:presupuesto}

La \autoref{tab:presupuesto} recoge el presupuesto de los componentes de la plataforma. La Raspberry Pi, el dongle LTE, la emisora con su receptor FlySky FS-i6X y el material auxiliar ya se tenían de antes de empezar el proyecto (no se compraron para él), por lo que su coste se estima a precio aproximado de reposición en vez de a partir de un ticket de compra real.

\begin{longtable}{p{0.34\textwidth}p{0.12\textwidth}p{0.18\textwidth}p{0.18\textwidth}}
\caption{Presupuesto inicial del sistema.}\label{tab:presupuesto}\\
\toprule
\textbf{Elemento} & \textbf{Cantidad} & \textbf{Coste unitario} & \textbf{Subtotal} \\
\midrule
\endfirsthead
\toprule
\textbf{Elemento} & \textbf{Cantidad} & \textbf{Coste unitario} & \textbf{Subtotal} \\
\midrule
\endhead
Chasis F450 con tren de aterrizaje & 1 & 16.29€ & 16.29€ \\
Motores brushless 2212 920KV & 4 & 6.50€ & 26.00€ \\
Hélices 1045 CW/CCW & 1 lote de 10 & 6.67€ & 6.67€ \\
Batería LiPo 4S 5200 mAh 60C & 1 & 31.19€ & 31.19€ \\
Stack H743 6S con ESC 55 A & 1 & 51.69€ & 51.69€ \\
Módulo GPS M10 con brújula & 1 & 14.28€ & 14.28€ \\
BEC ajustable & 1 & 4.39€ & 4.39€ \\
Placa antivibración & 1 & 4.99€ & 4.99€ \\
Cámara USB global shutter & 1 & 65.99€ & 65.99€ \\
Raspberry Pi (ya disponible, coste aprox. de reposición) & 1 & 60.00€ & 60.00€ \\
Dongle LTE (ya disponible, coste aprox. de reposición) & 1 & 20.00€ & 20.00€ \\
Emisora y receptor FlySky FS-i6X (ya disponible, coste aprox. de reposición) & 1 & 45.00€ & 45.00€ \\
Material auxiliar, conectores y fijaciones (coste aprox.) & 1 lote & 10.00€ & 10.00€ \\
\midrule
\textbf{Total} & & & \textbf{356.49€} \\
\bottomrule
\end{longtable}

% ---------------------------------------------------------------------